\begin{tomtat}
	\begin{minipage}[t]{0.3\textwidth}
		\vspace{0pt}
		\begin{tomtat}
			\begin{center}
				\textbf{ĐỀ 06}
			\end{center}
			
		\end{tomtat}
	\end{minipage}
	\hfill
	\begin{minipage}[t]{0.5\textwidth}
		\vspace{0pt}
		\begin{tomtat}
			\begin{center}
				\includegraphics[scale=0.08]{images/6}
			\end{center}
		\end{tomtat}
	\end{minipage}
\end{tomtat}
\textbf{I. PHẦN TRẮC NGHIỆM}
\begin{ex}
	Phương trình $(x-2)(2x-3)=0$ có nghiệm là ?
	\choice{$x=2$ và $x=3$}{$x=-2$ và $x=3/2$}{$x=2$ và $x=3/2$}{$x=-2$ và $x=-3/2$}
	\loigiai{
		Ta có phương trình tích: $(x-2)(2x-3)=0$.
		Điều này xảy ra khi một trong hai thừa số bằng 0:
		\begin{itemize}
			\item $x-2=0 \Rightarrow x=2$
			\item $2x-3=0 \Rightarrow 2x=3 \Rightarrow x=\frac{3}{2}$
		\end{itemize}
		Vậy nghiệm của phương trình là $x=2$ và $x=\frac{3}{2}$.
		\\ \textbf{ĐÁP ÁN C}
	}
\end{ex}

\begin{ex}
	Giá trị biểu thức $A=\frac{1}{2-\sqrt{3}}+\frac{1}{2+\sqrt{3}}$ là :
	\choice{1}{2}{3}{4}
	\loigiai{
		Ta quy đồng mẫu số:
		\begin{align*}
			A &= \frac{1(2+\sqrt{3})}{(2-\sqrt{3})(2+\sqrt{3})} + \frac{1(2-\sqrt{3})}{(2+\sqrt{3})(2-\sqrt{3})} \\
			&= \frac{2+\sqrt{3}}{2^2 - (\sqrt{3})^2} + \frac{2-\sqrt{3}}{2^2 - (\sqrt{3})^2} \\
			&= \frac{2+\sqrt{3}}{4 - 3} + \frac{2-\sqrt{3}}{4 - 3} \\
			&= \frac{2+\sqrt{3}}{1} + \frac{2-\sqrt{3}}{1} \\
			&= 2+\sqrt{3} + 2-\sqrt{3} \\
			&= 4
		\end{align*}
		\\ \textbf{ĐÁP ÁN D}
	}
\end{ex}

\begin{ex}
	Tìm giá trị của a, biết đồ thị hàm số $y=ax^2$ đi qua điểm M(-1;2).
	\choice{-1}{-2}{2}{D.-2} % Chú ý: Đáp án D có vẻ bị lặp lại và có thêm chữ D.
	\loigiai{
		Đồ thị hàm số $y=ax^2$ đi qua điểm M(-1;2) có nghĩa là khi $x=-1$ thì $y=2$.
		Thay $x=-1$ và $y=2$ vào phương trình hàm số, ta được:
		\begin{align*}
			2 &= a(-1)^2 \\
			2 &= a \cdot 1 \\
			a &= 2
		\end{align*}
		Vậy giá trị của $a$ là 2.
		\\ \textbf{ĐÁP ÁN C}
	}
\end{ex}

\begin{ex}
	Biểu thức $\sqrt[3]{27}-\sqrt{11-6\sqrt{2}}$ có giá trị bằng :
	\choice{$\sqrt{2}$}{$\sqrt{3}$}{2}{3}
	\loigiai{
		Ta có $\sqrt[3]{27} = 3$.
		Xét biểu thức $\sqrt{11-6\sqrt{2}}$:
		$11-6\sqrt{2} = 9 - 2 \cdot 3 \cdot \sqrt{2} + 2 = (3)^2 - 2 \cdot 3 \cdot \sqrt{2} + (\sqrt{2})^2 = (3-\sqrt{2})^2$.
		Vì $3 > \sqrt{2}$ (do $9 > 2$), nên $3-\sqrt{2} > 0$.
		Do đó, $\sqrt{11-6\sqrt{2}} = \sqrt{(3-\sqrt{2})^2} = |3-\sqrt{2}| = 3-\sqrt{2}$.
		Vậy, biểu thức đã cho bằng:
		$$ 3 - (3-\sqrt{2}) = 3 - 3 + \sqrt{2} = \sqrt{2} $$
		\\ \textbf{ĐÁP ÁN A}
	}
\end{ex}

\begin{ex}
	Cho tam giác ABC vuông tại A có AB = 3, AC=4. Khi đó số đo góc B bằng (làm tròn đến độ)?
	\choice{$52^\circ$}{$53^\circ$}{$54^\circ$}{$55^\circ$}
	\loigiai{
		Trong tam giác ABC vuông tại A:
		Ta có $\tan B = \frac{\text{cạnh đối}}{\text{cạnh kề}} = \frac{AC}{AB} = \frac{4}{3}$.
		Để tìm số đo góc B, ta sử dụng máy tính cầm tay:
		$B = \arctan\left(\frac{4}{3}\right) \approx 53.13^\circ$.
		Làm tròn đến độ, ta được $B \approx 53^\circ$.
		\\ \textbf{ĐÁP ÁN B}
	}
\end{ex}

\begin{ex}
	Tính thể tích V của hình cầu có bán kính $R = 3\text{cm}$?.
	\choice{$V=12\pi(\text{cm}^3)$}{$V=18\pi(\text{cm}^3)$}{$V=27\pi(\text{cm}^3)$}{$V=36\pi(\text{cm}^3)$}
	\loigiai{
		Thể tích của hình cầu được tính bằng công thức: $V = \frac{4}{3}\pi R^3$.
		Với $R=3 \text{ cm}$:
		\begin{align*}
			V &= \frac{4}{3}\pi (3)^3 \\
			&= \frac{4}{3}\pi \cdot 27 \\
			&= 4\pi \cdot \frac{27}{3} \\
			&= 4\pi \cdot 9 \\
			&= 36\pi (\text{cm}^3)
		\end{align*}
		\\ \textbf{ĐÁP ÁN D}
	}
\end{ex}

\begin{ex}
	Hình vành khuyên giới hạn bởi hai đường tròn $(O; 2 \text{ cm})$ và $(O; 4 \text{ cm})$ có diện tích bằng bao nhiêu?
	\choice{$S=12\pi(\text{cm}^2)$}{$S=8\pi(\text{cm}^2)$}{$S=6\pi(\text{cm}^2)$}{$S=4\pi(\text{cm}^2)$}
	\loigiai{
		Diện tích hình tròn lớn (bán kính $R_L = 4 \text{ cm}$) là $S_L = \pi R_L^2 = \pi (4)^2 = 16\pi \text{ cm}^2$.
		Diện tích hình tròn nhỏ (bán kính $R_N = 2 \text{ cm}$) là $S_N = \pi R_N^2 = \pi (2)^2 = 4\pi \text{ cm}^2$.
		Diện tích hình vành khuyên là hiệu diện tích của hình tròn lớn và hình tròn nhỏ:
		$$ S = S_L - S_N = 16\pi - 4\pi = 12\pi (\text{cm}^2) $$
		\\ \textbf{ĐÁP ÁN A}
	}
\end{ex}

\begin{ex}
	Ba bạn An, Bình, Hà được xếp ngẫu nhiên ngồi trên một hàng ghế có ba chỗ ngồi. Xác suất của biến cố: An không ngồi ngoài cùng bên phải là:
	\choice{1/2}{1/3}{2/3}{1/4}
	\loigiai{
		Gọi 3 vị trí ngồi là 1 (trái), 2 (giữa), 3 (phải).
		Tập hợp tất cả các kết quả có thể xảy ra khi xếp 3 bạn An (A), Bình (B), Hà (H) vào 3 chỗ ngồi (không gian mẫu $\Omega$):
		\begin{itemize}
			\item (A, B, H)
			\item (A, H, B)
			\item (B, A, H)
			\item (B, H, A)
			\item (H, A, B)
			\item (H, B, A)
		\end{itemize}
		Số phần tử của không gian mẫu là $n(\Omega) = 6$.
		
		Tập hợp các kết quả thuận lợi cho biến cố "An không ngồi ngoài cùng bên phải":
		Điều này có nghĩa là An ngồi ở vị trí 1 (trái) hoặc vị trí 2 (giữa).
		\begin{itemize}
			\item Nếu An ngồi vị trí 1 (trái): (A, B, H), (A, H, B) - có 2 trường hợp.
			\item Nếu An ngồi vị trí 2 (giữa): (B, A, H), (H, A, B) - có 2 trường hợp.
		\end{itemize}
		Vậy, các kết quả thuận lợi cho biến cố là: (A, B, H), (A, H, B), (B, A, H), (H, A, B).
		Số phần tử của biến cố thuận lợi là $n(A) = 4$.
		
		Xác suất của biến cố là:
		$$ P(A) = \frac{n(A)}{n(\Omega)} = \frac{4}{6} = \frac{2}{3} $$
		\\ \textbf{ĐÁP ÁN C}
	}
\end{ex}
\newpage
\textbf{II. PHẦN TỰ LUẬN}
\begin{ex}
	Rút gọn biểu thức $B = \left(\frac{1}{\sqrt{x}} + \frac{1}{3-\sqrt{x}}\right) : \frac{1}{3-\sqrt{x}}$ với $x > 0, x \ne 9$.
	\loigiai{
		Với $x > 0, x \ne 9$, ta có:
		\begin{align*}
			B &= \left(\frac{1}{\sqrt{x}} + \frac{1}{3-\sqrt{x}}\right) : \frac{1}{3-\sqrt{x}} \\
			&= \left(\frac{3-\sqrt{x} + \sqrt{x}}{\sqrt{x}(3-\sqrt{x})}\right) \cdot (3-\sqrt{x}) \\
			&= \frac{3}{\sqrt{x}(3-\sqrt{x})} \cdot (3-\sqrt{x})
		\end{align*}
		Vì $x \ne 9$ nên $3-\sqrt{x} \ne 0$. Ta có thể rút gọn $3-\sqrt{x}$.
		\begin{align*}
			B &= \frac{3}{\sqrt{x}}
		\end{align*}
		Vậy $B = \frac{3}{\sqrt{x}}$.
	}
\end{ex}

\begin{ex}
	Giải hệ phương trình sau: $\begin{cases} 2x-3y=7 \\ x+5y=-3 \end{cases}$.
	\loigiai{
		Ta có hệ phương trình:
		$$ \begin{cases} 2x-3y=7 & (1) \\ x+5y=-3 & (2) \end{cases} $$
		Từ phương trình (2), suy ra $x = -3 - 5y$.
		Thay $x = -3 - 5y$ vào phương trình (1):
		\begin{align*}
			2(-3-5y) - 3y &= 7 \\
			-6 - 10y - 3y &= 7 \\
			-13y &= 7 + 6 \\
			-13y &= 13 \\
			y &= -1
		\end{align*}
		Thay $y=-1$ vào $x = -3 - 5y$:
		\begin{align*}
			x &= -3 - 5(-1) \\
			x &= -3 + 5 \\
			x &= 2
		\end{align*}
		Vậy nghiệm của hệ phương trình là $(x; y) = (2; -1)$.
	}
\end{ex}

\begin{ex}
	Cho phương trình $x^2 - 2x - 5 = 0$. Không giải phương trình hãy tính giá trị biểu thức $A = \frac{x_1^2}{x_2} + \frac{x_2^2}{x_1} - 5x_1 - 5x_2$.
	\loigiai{
		Phương trình đã cho là $x^2 - 2x - 5 = 0$.
		Đây là phương trình bậc hai có dạng $ax^2+bx+c=0$ với $a=1, b=-2, c=-5$.
		Ta tính biệt thức $\Delta = b^2-4ac = (-2)^2 - 4(1)(-5) = 4 + 20 = 24$.
		Vì $\Delta = 24 > 0$ nên phương trình có hai nghiệm phân biệt $x_1, x_2$ (và $x_1, x_2 \ne 0$ vì $c \ne 0$).
		
		Theo hệ thức Vi-ét, ta có:
		$$ \begin{cases} 
			x_1+x_2 = \dfrac{-b}{a} = \dfrac{-(-2)}{1} = 2 \\ 
			x_1x_2 = \dfrac{c}{a} = \dfrac{-5}{1} = -5 
		\end{cases} $$
		Biến đổi biểu thức A:
		\begin{align*}
			A &= \frac{x_1^3 + x_2^3}{x_1x_2} - 5(x_1+x_2) \\
			&= \frac{(x_1+x_2)(x_1^2 - x_1x_2 + x_2^2)}{x_1x_2} - 5(x_1+x_2) \\
			&= \frac{(x_1+x_2)((x_1+x_2)^2 - 2x_1x_2 - x_1x_2)}{x_1x_2} - 5(x_1+x_2) \\
			&= \frac{(x_1+x_2)((x_1+x_2)^2 - 3x_1x_2)}{x_1x_2} - 5(x_1+x_2)
		\end{align*}
		Thay giá trị $x_1+x_2 = 2$ và $x_1x_2 = -5$ vào biểu thức $A$:
		\begin{align*}
			A &= \frac{(2)((2)^2 - 3(-5))}{-5} - 5(2) \\
			&= \frac{2(4 + 15)}{-5} - 10 \\
			&= \frac{2 \cdot 19}{-5} - 10 \\
			&= \frac{38}{-5} - 10 \\
			&= -\frac{38}{5} - \frac{50}{5} \\
			&= -\frac{88}{5}
		\end{align*}
		Vậy giá trị của biểu thức $A$ là $-\frac{88}{5}$.
	}
\end{ex}

\begin{ex}
	Hai ô tô cùng xuất phát tại A đi trên cùng quãng đường từ A đến B dài 200 km. Biết mỗi giờ xe thứ nhất chạy nhanh hơn xe thứ 2 là 10 km. Nên xe thứ nhất đến B trước xe thứ 2 là 1 giờ. Tính Tốc độ mỗi xe?
	\loigiai{
		Gọi tốc độ của xe thứ hai là $x$ (km/h). Điều kiện $x > 0$.
		Thì tốc độ của xe thứ nhất là $x+10$ (km/h).
		
		Thời gian xe thứ hai đi hết quãng đường AB là: $ \frac{200}{x}$ (giờ).
		Thời gian xe thứ nhất đi hết quãng đường AB là: $\frac{200}{x+10}$ (giờ).
		
		Vì xe thứ nhất đến B trước xe thứ hai 1 giờ, nên ta có phương trình:
		$$ \frac{200}{x} - \frac{200}{x+10} = 1 $$
		Quy đồng mẫu số:
		\begin{align*}
			\frac{200(x+10) - 200x}{x(x+10)} &= \frac{x(x+10)}{x(x+10)} \\
			200x + 2000 - 200x &= x^2 + 10x \\
			2000 &= x^2 + 10x \\
			x^2 + 10x - 2000 &= 0
		\end{align*}
		Tính biệt thức $\Delta = b^2 - 4ac = 10^2 - 4(1)(-2000) = 100 + 8000 = 8100$.
		$\sqrt{\Delta} = \sqrt{8100} = 90$.
		Phương trình có hai nghiệm:
		$$ x_1 = \frac{-10 + 90}{2(1)} = \frac{80}{2} = 40 $$
		$$ x_2 = \frac{-10 - 90}{2(1)} = \frac{-100}{2} = -50 $$
		Vì $x > 0$, nên ta nhận $x = 40$.
		Vậy, tốc độ của xe thứ hai là 40 km/h.
		Tốc độ của xe thứ nhất là $x+10 = 40+10 = 50$ km/h.
		
		Kiểm tra lại:
		Thời gian xe thứ hai đi: $\frac{200}{40} = 5$ giờ.
		Thời gian xe thứ nhất đi: $\frac{200}{50} = 4$ giờ.
		Xe thứ nhất đến trước xe thứ hai: $5 - 4 = 1$ giờ (đúng).
		
		Vậy tốc độ xe thứ nhất là 50 km/h và tốc độ xe thứ hai là 40 km/h.
	}
\end{ex}

\begin{ex}
	Sau buổi kiểm tra môn Toán, giáo viên ghi lại số điểm mà học sinh đạt được như sau:
	\begin{center}
	\begin{tabular}{|c|c|c|c|c|c|c|c|c|c|}
		\hline
		5 & 6 & 8 & 7 & 4 & 5 & 7 & 6 & 8 & 9 \\
		\hline
		7 & 5 & 6 & 8 & 7 & 8 & 8 & 9 & 8 & 7 \\
		\hline
		5 & 6 & 5 & 6 & 5 & 4 & 5 & 6 & 10 & 7 \\
		\hline
		6 & 5 & 7 & 6 & 5 & 8 & 5 & 5 & 6 & 5 \\
		\hline
	\end{tabular}
\end{center}
	Hỏi mẫu số liệu trên gồm những giá trị khác nhau nào? Hãy lập bảng tần số và bảng tần số tương đối của số điểm mà học sinh đạt được.
	\loigiai{
		Các giá trị khác nhau của điểm số trong mẫu số liệu là: 4, 5, 6, 7, 8, 9, 10.
		Tổng số học sinh (kích thước mẫu) là $N = 4 \times 10 = 40$.
		\begin{center}
			\begin{tabular}{|l|c|c|c|c|c|c|c|c|}
				\hline
				Điểm số (Giá trị $x_i$) & 4 & 5 & 6 & 7 & 8 & 9 & 10 & \textbf{Tổng} \\
				\hline
				Tần số ($n_i$) & 2 & 11 & 9 & 8 & 7 & 2 & 1 & \textbf{N = 40} \\
				\hline
				Tần số tương đối ($f_i$) & $5\%$ & $27,5\%$ & $22,5\%$ & $20\%$ & $17,5\%$ & $5\%$ & $2,5\%$ & \textbf{100\%} \\
				\hline
			\end{tabular}
		\end{center}
	}
\end{ex}
\begin{ex}
	Cho $\triangle ABC$ nhọn, đường tròn $(O)$ đường kính $BC$ cắt $AB, AC$ lần lượt tại $E$ và $F$. Gọi $H$ là giao điểm của $BF$ và $CE$, $AH$ cắt $BC$ tại $D$.
	\begin{enumerate}
		\item[a)] Chứng minh: Tứ giác $AEHF$ nội tiếp, xác định tâm $K$ của đường tròn này.
		\item[b)] Chứng minh: $KE$ là tiếp tuyến của đường tròn $(O)$ và năm điểm $O, D, E, K, F$ cùng thuộc một đường tròn.
	\end{enumerate}	
\loigiai{
		\begin{center}
			\begin{tikzpicture}[line join = round, line cap=round,>=stealth,font=\footnotesize,scale=1]
				% Def
				\def\r{3} % r=3cm
				
				% Coordinates
				\coordinate (O) at (0,0);
				\coordinate (B) at (-\r,0);
				\coordinate (C) at (\r,0);
				\coordinate (A) at (110:1.6*\r); % A(110 độ, 2.1*r)
				
				\coordinate (E) at ($(A)!(C)!(B)$); % E la hinh chieu cua C len AB
				\coordinate (F) at ($(A)!(B)!(C)$); % F la hinh chieu cua B len AC
				\coordinate (D) at ($(B)!(A)!(C)$); % D la hinh chieu cua A len BC
				
				% De tim H, can \usetikzlibrary{intersections} o preamble
				\path[name path=CE_path] (C) -- (E);
				\path[name path=BF_path] (B) -- (F);
				\coordinate[name intersections={of=CE_path and BF_path, by=H}]; % H la giao diem CE va BF
				
				\coordinate (K) at ($(A)!0.5!(H)$); % K la trung diem cua AH
				
				% Draw
				\draw (O) circle (\r); % Ve duong tron (O,r)
				\draw (A) -- (B) -- (C) -- cycle; % Ve tam giac ABC
				\draw (A) -- (D); % Ve AD
				\draw (B) -- (F); % Ve BF
				\draw (C) -- (E); % Ve CE
				
				% Ky hieu goc vuong bang lenh \pic
				% Can \usetikzlibrary{angles, quotes} o preamble
				% Goc AEH (thuc ra la AEC hoac CEB vuong)
				\pic[draw, angle radius=2.5mm]{right angle = C--E--B}; % E la chan duong cao tu C xuong AB nen CEB vuong
				% Goc HFA (thuc ra la BFC hoac CFA vuong)
				\pic[draw, angle radius=2.5mm]{right angle = B--F--C}; % F la chan duong cao tu B xuong AC nen BFC vuong
				% Goc ADB vuong
				\pic[draw, angle radius=2.5mm]{right angle = A--D--B}; 
				
				% Fill 
				\fill (O) circle (1.5pt);
				\fill (A) circle (1.5pt);
				\fill (B) circle (1.5pt);
				\fill (C) circle (1.5pt);
				\fill (E) circle (1.5pt);
				\fill (F) circle (1.5pt);
				\fill (D) circle (1.5pt);
				\fill (H) circle (1.5pt);
				\fill (K) circle (1.5pt);
				
				% Nodes	
				\node[below left] at (O) {$O$};
				\node[above] at (A) {$A$};
				\node[below left] at (B) {$B$};
				\node[below right] at (C) {$C$};
				\node[above left] at (E) {$E$}; % Vi tri nhan co the can dieu chinh
				\node[above right] at (F) {$F$}; % Vi tri nhan co the can dieu chinh
				\node[below] at (D) {$D$}; % Vi tri nhan co the can dieu chinh
				\node[above] at (H) {$H$}; % Vi tri nhan co the can dieu chinh
				\node[above right] at (K) {$K$}; % Vi tri nhan co the can dieu chinh
				% Ve va danh dau doan thẳng bằng nhau
			\end{tikzpicture}
			
		\end{center}
			\begin{enumerate}
			\item Chứng minh: $KE$ là tiếp tuyến của đường tròn $(O)$
			
			$\widehat{BEC} = \widehat{BFC} = 90^\circ$ (Góc nội tiếp chắn nửa đường tròn)
			
			Suy ra BF, CE là 2 đường cao của $\triangle ABC$. Suy ra H là trực tâm của $\triangle ABC$. Suy ra $AH \perp BC$.
			
			Ta có: $\widehat{AEK} = \widehat{KAE}$ ($\triangle AKE$ cân tại K)
			
			$\widehat{BEO} = \widehat{EBO}$ ($\triangle BOE$ cân tại O). Suy ra: $\widehat{AEK} + \widehat{BEO} = \widehat{KAE} + \widehat{EBO}$
			
			Mà $\widehat{KAE} + \widehat{EBO} = 90^\circ$ ($\triangle ADB$ vuông tại D). Nên $\widehat{KEO} = 90^\circ$.
			
			Suy ra KE là tiếp tuyến của đường tròn (O).
			\item năm điểm $O, D, E, K, F$ cùng thuộc một đường tròn.
			
			Chứng minh tương tự: KF là tiếp tuyến của (O).
			
			Vì $\widehat{KEO} = \widehat{KDO} = \widehat{KFO} = 90^\circ$ nên KO là đường kính của đường tròn ngoại tiếp các tam giác vuông KEO, KDO, KFO. Vậy năm điểm O, D, E, K, F cùng thuộc một đường tròn đường kính OK.
			
		\end{enumerate}
		}
\end{ex}

\begin{ex}
	\begin{enumerate}
		\item[a)] Tính chiều cao của một cột tháp, biết rằng khi các tia nắng mặt trời tạo với mặt đất một góc bằng $50^\circ$ thì bóng của nó trên mặt đất dài $96 \text{ m}$. (Kết quả làm tròn đến chữ số thập phân thứ nhất).
		\item[b)] Cho $a,b,c$ thực dương thỏa mãn $abc = 1$. Chứng minh rằng:
		$$ \frac{1}{\sqrt{a^4-a^3+ab+2}} + \frac{1}{\sqrt{b^4-b^3+bc+2}} + \frac{1}{\sqrt{c^4-c^3+ac+2}} \le \sqrt{3} $$
	\end{enumerate}
	\loigiai{
		\begin{enumerate}
		\item Chiều cao của cột tháp là : $96 \cdot \tan50^\circ = 11,4 (m)$
		\item Ta có: $(a-1)^2(a^2+a+1) \ge 0 \Leftrightarrow (a^2-2a+1)(a^2+a+1) \ge 0$
		
		$a^4-a^3-a+1 \ge 0$
		
		$a^4-a^3+1 \ge a \Leftrightarrow a^4-a^3+ab+2 \ge ab+a+1$
		
		$\frac{1}{\sqrt{a^4-a^3+ab+2}} \le \frac{1}{\sqrt{ab+a+1}}$
		Chứng minh hoàn toàn tương tự ta có:
		
		$\frac{1}{\sqrt{b^4-b^3+bc+2}} \le \frac{1}{\sqrt{bc+b+1}}$; $\frac{1}{\sqrt{c^4-c^3+ac+2}} \le \frac{1}{\sqrt{ac+c+1}}$
		
		Như vậy:
		
		$VT \le \frac{1}{\sqrt{ab+a+1}} + \frac{1}{\sqrt{bc+b+1}} + \frac{1}{\sqrt{ac+c+1}} \le \sqrt{3(\frac{1}{ab+a+1} + \frac{1}{bc+b+1} + \frac{1}{ac+c+1})}$
		(Áp dụng BĐT Bunyakovski cho 3 số)
		
		Lại có
		$\sqrt{3(\frac{1}{ab+a+1} + \frac{1}{bc+b+1} + \frac{1}{ac+c+1})} = \sqrt{3(\frac{1}{ab+a+1} + \frac{a}{abc+ab+a} + \frac{ab}{a^2bc+abc+ab})}$
		$= \sqrt{3(\frac{1}{ab+a+1} + \frac{a}{1+ab+a} + \frac{ab}{a+ab+1})} = \sqrt{3}$
		
		Vậy ta có điều phải chứng minh
		
		Dấu “=” xảy ra khi và chỉ khi $a=b=c=1$.
		
	\end{enumerate}
	}
\end{ex}
\setcounterref{ex}{0}